\section{Experiments}
\label{sec:experiments}

% TARGET: ~2.5 pages

% PLACEHOLDER — to be written

\subsection{Setup}
\label{sec:setup}
% Datasets: eICU (source), MIMIC-IV (target), HiRID (2nd source)
% Tasks: Mortality24 (per-stay), AKI (per-timestep), Sepsis (per-timestep),
%        LoS (per-timestep, regression), KidneyFunction (per-stay, regression)
% Frozen model: LSTM trained on MIMIC-IV via YAIB benchmark
% Baselines: 8 DA methods (DANN, CORAL, CoDATS, CDAN, RAINCOAT, CLUDA, fine-tuning, statistics-only)
% Metrics: AUROC, AUCPR for classification; MAE, MSE for regression
% Implementation: PyTorch, batch_size=16, early stopping, seeds

\subsection{Main Results}
\label{sec:main_results}
% Table 1: All tasks x all methods (ours + 8 baselines + frozen baseline + native models)
% Key findings:
%   - Surpasses eICU-native LSTM on all 3 classification tasks
%   - AKI: surpasses MIMIC-native (+0.056 AUROC)
%   - Outperforms all 8 DA baselines by 2-4x

\subsection{Multi-Source Generalization}
\label{sec:multi_source}
% Table 2: HiRID -> MIMIC results
% Even larger gains: AKI +0.078, Sepsis +0.078, Mortality +0.047

\subsection{Architecture-Agnostic Transfer}
\label{sec:mas}
% Table 3: LSTM-trained translator evaluated on frozen GRU, TCN
% No retraining needed; 47-86% of LSTM gain preserved

\subsection{Ablation Studies}
\label{sec:ablations}
% Table 4: Component ablation (one variable at a time)
% Key ablations: fidelity loss, cross-attention layers, memory bank, feature gate
% n_cross_layers analysis: 3 for dense labels, 2 for sparse

% \subsection{Analysis}
% \label{sec:analysis}
% Optional: t-SNE/UMAP visualization, training dynamics, gradient analysis
